% Source: https://artofproblemsolving.com/community/c3h1399580_picks_theorem
% Stable identity: aops-topic-1399580-q01
Let the vertices be $A=(x_1,y_1)$, $B=(x_2,y_2)$, $C=(x_3,y_3)$ with integer coordinates. By the shoelace formula,
\[
[ABC]=\dfrac{1}{2}\bigl|x_1(y_2-y_3)+x_2(y_3-y_1)+x_3(y_1-y_2)\bigr| .
\]
The expression inside the absolute value is an integer $K$, so $[ABC]=\dfrac{K}{2}$ with $K$ a nonnegative integer: the area is a multiple of $\dfrac{1}{2}$.

\begin{asy}
usepackage("amsmath");
size(5cm);
defaultpen(fontsize(10pt));
for(int i=0;i<=6;++i) for(int j=0;j<=4;++j) dot((i,j),gray(0.6)+2);
pair A=(0,0),B=(5,1),C=(2,4);
fill(A--B--C--cycle,paleblue); draw(A--(5,0)--(5,4)--(0,4)--cycle,dashed+gray);
for(int i=0;i<=6;++i) for(int j=0;j<=4;++j) dot((i,j),gray(0.6)+2);
draw(A--B--C--cycle,blue+linewidth(1));
dot(A);dot(B);dot(C);
label("$A$",A,SW); label("$B$",B,E); label("$C$",C,N);
label("\mbox{}",(min(currentpicture,user=true)+max(currentpicture,user=true))/2,black); // keeps the text after the figure black when LaTeX typesets the labels inline
\end{asy}

(For instance, the triangle above has area $\dfrac{1}{2}\left|0(1-4)+5(4-0)+2(0-1)\right|=\dfrac{18}{2}=9$; equivalently it is the $5\times4$ rectangle minus three right triangles of areas $\dfrac{5}{2}$, $\dfrac{9}{2}$ and $4$.)

Now $1223442.25=1223442+\dfrac{1}{4}$ is not a multiple of $\dfrac{1}{2}$, so no lattice triangle has this area, and the student's area computation contains an error. Pick's theorem $A=I+\dfrac{B}{2}-1$ is consistent with this: for any lattice polygon $I$ and $B$ are integers, so $A$ must be a multiple of $\dfrac{1}{2}$, and a non-integer value of $I$ is a sign that the area was wrong, not that Pick's theorem fails.
